% 對應英文：sections/03_related_work.tex

\section{橫截面報酬預測的機器學習方法}

\citet{gu2020empirical} 為機器學習橫截面資產定價提供了經典基礎，
指出在特徵集豐富的情況下，
樹模型與神經網路均能超越傳統線性模型，
而真正的預測力往往來自特徵之間的非線性交互作用，而非單一指標本身。
\citet{freyberger2020dissecting} 亦從非參數角度證實，
當資產特徵之間存在複雜交互項時，
線性模型對平均報酬的近似能力有限。
這組文獻的重要啟示在於：
若橫截面特徵已經足夠豐富，模型的關鍵任務不一定是建構更長的時序記憶，
而是學會如何比較當期不同資產之間的相對位置。

在加密貨幣市場中，\citet{liu2022cryptocurrency} 建立了共同風險因子的基本框架，
顯示動量與規模等效果在加密資產中確實存在，但其形成方式與股票市場不同。
更直接相關的是 \citet{cakici2024crypto_ml} 與 \citet{filippou2024crypto_ml}：
前者以多種機器學習方法研究加密貨幣橫截面報酬，
指出報酬可預測性大致來自少數簡單特徵，模型複雜度的邊際報酬有限；
後者則發現，在納入網路活動、動量、技術指標與投資人關注等資料後，
非線性模型可產生顯著的樣本外預測力與經濟價值。
這些研究說明，加密貨幣市場的可預測性並非不存在，
而是取決於模型如何對異質特徵進行整合與排序。
本文的 CS-Gated 架構，正是沿著這條脈絡，
直接測試「僅依賴當期橫截面訊號」是否已足以生成高品質排序。

\section{金融領域的學習排序方法}

學習排序文獻 \citep{cao2007learning} 為橫截面投資組合建構提供了自然的理論基礎。
與直接預測報酬值不同，排序方法的目標是讓高分資產在橫截面上穩定地排在低分資產之前，
因此更貼近多空投資組合的實務需求。
\citet{poh2021building} 是與本文最直接相關的先驅工作：
他們指出，若直接將原始報酬放入 ListNet 的 softmax 目標，
會因不同期別的報酬尺度差異而造成訓練不穩定；
改用橫截面排名百分位後，模型在股票橫截面上的排序績效顯著改善。

本文完整承襲其核心想法，但將其置於波動更高的加密貨幣市場中重新檢驗。
相較於股票市場，加密貨幣之間的波動率差距更大，
同一週內不同資產的報酬分布也更容易出現厚尾與極端值。
因此，目標函數尺度不一致在本研究設定下更加嚴重，
也使排名正規化更具有方法論上的必要性。
此外，本文再加入梯度累積，以改善每一步 ListNet 梯度所依賴的橫截面經驗分布品質，
藉此將排序方法從「可訓練」進一步推向「穩定可重現」。

\section{時序深度學習架構}

LSTM \citep{hochreiter1997long,fischer2018lstm_finance} 長期被視為金融時間序列中的穩健基準，
其優勢在於能以相對節制的參數量整合有限回看窗口中的局部時序依賴。
對於社群媒體發文量、ETF 資金流或其他具有事件延續性的訊號而言，
LSTM 能有效捕捉短期序列模式，
因此仍是本文不可或缺的比較基準。

另一方面，TFT \citep{lim2021temporal} 結合變數選擇網路、LSTM 編碼器與多頭注意力，
強調對異質輸入的可解釋整合。
相較一般 Transformer 架構 \citep{vaswani2017attention}，
TFT 更適合處理同時包含靜態與時變特徵的金融場景。
然而，TFT 的優勢通常伴隨較高模型容量與較高訓練不穩定性，
因此本文特別修正其變數選擇網路，使其更接近原始規格，
並與 LSTM 及純橫截面架構進行並列比較。
這樣的設計使我們得以回答一個更具研究價值的問題：
究竟是「時序編碼能力」本身帶來改善，
還是「正確的損失函數與特徵配置」才是績效差異的主要來源。

\section{替代資料來源}

替代資料在金融預測中的應用已相當廣泛。
\citet{tetlock2007giving} 證明媒體內容可影響股票短期報酬，
\citet{bollen2011twitter} 則顯示 Twitter 心情指標能在一定程度上領先市場變動。
在加密貨幣領域，\citet{maitre2025social} 進一步指出，
社群媒體注意力與加密貨幣橫截面報酬存在顯著關聯，
代表文字與敘事熱度已成為不可忽視的市場資訊來源。

本文使用的川普社群媒體特徵與上述文獻不同，
並非衡量整體市場情緒，而是捕捉單一高影響力政策人物的訊號。
這種設計的意義在於將「政策敘事」視為可量化的結構性衝擊來源，
而非單純的雜訊情緒代理變數。
同時，ETF 資金流資料則補上鏈上資料無法直接觀察到的機構需求面。
自 2024 年 1 月現貨比特幣 ETF 核准後，
\citet{bitcoin_etf_farside} 所提供的細粒度流量資料，
使研究者得以區分 Grayscale 存量資金贖回與新增資金流入，
這對於理解事件後的供需結構具有獨特價值。
本文將這兩類資料與價格、技術、鏈上特徵放入同一排序框架，
目的不只是追求更高績效，
更是希望辨識不同資料型態對不同模型架構的適配性。

\section{統計推斷方法}

金融機器學習研究若只報告樣本外績效而不處理序列相關性與非常態問題，
容易高估策略的統計可信度。
因此，本文採用 Newey--West
\citep{newey1987simple,newey1994automatic}
異方差與自相關一致（HAC）$t$ 統計量，
並使用 Politis 與 Romano \citep{politis1994stationary}
提出的定態區塊自助法來構建夏普比率的信賴區間。
這些方法並不直接提高模型績效，
但可避免將短樣本、集中的市場行情或偶然極端報酬誤判為穩健結論，
對於僅有 49 週測試期的本文而言尤其重要。
